\documentclass[12pt]{article}
\usepackage[utf8]{inputenc}
\usepackage[T1]{fontenc}
\usepackage{hyperref}
\title{Koguryo Language Resources for K-12 Curriculum Development}
\author{Compiled by Claude Sonnet 5.5 (claude-sonnet-5-5) for Andr\'as Kornai}
\date{}
\begin{document}
\maketitle

\section{Introduction}
Koguryo (Goguryeo, ISO 639-3 zkg) is an extinct language of the Goguryeo kingdom. The surviving record is small and fragmentary, and scholars disagree about its affiliation. Some link it to Japonic, some to Koreanic, some to Tungusic, and some say the evidence is insufficient. No native speakers, school materials or curricula were found in the searches. A K-12 curriculum that uses the language itself is therefore not feasible. Only history and linguistics enrichment about the language can be supported.

\section{Primary Educational Resources}
No primers, textbooks or school programmes for Koguryo were found.

\section{Linguistic and Reference Materials}
\begin{itemize}
\item Wikipedia overview of the Goguryeo language: \url{https://en.wikipedia.org/wiki/Goguryeo_language}
\item Wikipedia on Peninsular Japonic, which discusses the proposed link to Goguryeo: \url{https://en.wikipedia.org/wiki/Peninsular_Japonic}
\item Wikipedia on the placename glosses in the Samguk sagi, the main source of Goguryeo vocabulary: \url{https://demo.azizisearch.com/lite/wikipedia/page/Placename_glosses_in_the_Samguk_sagi}
\item A Korean-language academic article on reconstructing Goguryeo grammatical forms: \url{https://scholar.kyobobook.co.kr/article/detail/4010023444916}
\item A Korean-language article on Goguryeo vocabulary: \url{https://www.earticle.net/Article/A77784}
\end{itemize}

\section{Religious and Community Resources}
No religious or community-language resources were found.

\section{Digital and Audio Resources}
No audio, video or software resources were found.

\section{Curriculum Development Context}
The data are mostly a small set of words, about 80 to 100 from placenames and 139 entries in one Korean study. These could support a short unit on historical linguistics and the Goguryeo kingdom. They cannot support a language-medium curriculum.
\end{document}